\begin{center}
    {\large \textbf{Giorno 16-17} \textbar\ 40 ore sveglio \textbar\ 24-25 Agosto 2024 \textbar\ \textbf{Torino}, Italia}
\end{center}

\noindent\rule{\textwidth}{0.4pt}

\landscapeLargeMargin{images/day16/PXL_20240824_001901905.MP.jpg}

\landscapeLargeMargin{images/day16/PXL_20240824_035948160.MP.jpg}

\landscapeLargeMargin{images/day16/PXL_20240824_215722227.MP.jpg}

\landscapeLargeMargin{images/day16/PXL_20240824_215848306.jpg}

\landscapeLargeMargin{images/day16/PXL_20240824_231514589.jpg}

\landscapeLargeMargin{images/day16/PXL_20240825_015143099.jpg}

\landscapeLargeMargin{images/day16/PXL_20240825_052525291.jpg}

\noindent 24-25 Agosto 2024 \newline

Notte di riposo completa, che ci sarà senz'altro utile: ci aspetta un viaggio di quaranta ore. Il programma prevede un volo da Bali a Jakarta, un altro da Jakarta a Jeddah, un'attesa di dieci ore in aeroporto e infine un volo da Jeddah a Milano, dove un autobus ci porterà a casa.

Il giorno prima avevamo prenotato un Grab per le nove, così abbiamo tutto il tempo per fare colazione. Il buffet dell'albergo lascia molto a desiderare: c'è una vasta scelta di prodotti di ogni tipo e provenienza, ma la qualità è pessima. Non so perché, ma penso sempre agli americani, che cercano l'abbondanza senza saper apprezzare il cibo — e infatti li vedo con i piatti pieni di bacon e uova. Noi optiamo per una colazione leggera, dolci e omelette, godendoci almeno il pasto seduti a bordo piscina.

Finita la colazione facciamo rapidamente il check-out: il nostro Grab è in anticipo e vogliamo approfittarne. L'autista, di poche parole, è il più lento che ci sia mai capitato e passa tutto il tempo al cellulare, prima su Facebook e poi a leggere le notizie, mentre davanti a lui il traffico scorre. Siamo molto in anticipo quindi non ne faccio un problema, ma mi dà un po' fastidio che mi addebiti il costo del parcheggio in aeroporto tramite l'app senza dirmelo. La procedura è corretta, ma avrebbe potuto avvisarci.

Arrivati in aeroporto, proviamo subito a fare il check-in, ma ci informano che siamo troppo in anticipo e che non aprirà prima di un'ora. La modalità Marines "sbrigati e aspetta" è ufficialmente attivata. Passiamo l'oretta seduti a scrivere e rileggere alcune pagine del diario, e nel frattempo ricevo una notifica che mi avvisa di un cambiamento d'orario del volo Saudia da Jakarta. Per un attimo mi preoccupo, perché la seconda coincidenza non è collegata alla prima e la compagnia non è tenuta a garantire il tempo tra i due voli, ma poi scopro che il volo è solo posticipato di tre quarti d'ora. Meglio così: l'attesa di dieci ore a Jeddah sarà un po' più breve.

Quando finalmente arriva l'ora procediamo senza problemi al check-in. Il primo volo è operato da Air Asia, una compagnia con cui sappiamo di poter portare senza problemi la scatola dei souvenir in terracotta, e così è effettivamente. Passiamo i controlli e ci imbarchiamo senza troppe difficoltà, con un volo breve e tranquillo.

\begin{center}
    % Decorative line
    \noindent\rule{0.6\textwidth}{0.4pt}
\end{center}

Una volta arrivati a Jakarta, però, cominciano i problemi. All'andata non ce ne eravamo accorti, ma l'aeroporto è enorme e per raggiungere il Terminal 3, dedicato alle partenze internazionali, dobbiamo prendere lo SkyTrain. Arrivati al terminal, ci mettiamo in coda per imbarcare i bagagli e — probabilmente — scegliamo la fila sbagliata. Sbagliata non per nostra colpa, forse il termine più corretto è "sfigata", poiché era l'unica disponibile per chi aveva già fatto il check-in online, ma l'addetta ai terminali impiega letteralmente dieci minuti per ogni gruppo di passeggeri. Vediamo le file accanto alle nostre procedere spedite, servite da addetti molto più veloci, e perdiamo così più di un'ora in piedi in coda.

Quando è quasi il nostro turno, una delle addette delle file Priority ci fa cenno di avvicinarci e noi ci precipitiamo da lei prima ancora che abbia finito di alzare la mano. In pochi minuti ci fa i biglietti e imbarca i bagagli, ma ci informa che non sarà possibile imbarcare la scatola dei souvenir così com'è — sarà necessario reimpacchettarla in un formato diverso. Eravamo preparati a questa eventualità e tiriamo fuori una sacca morbida da trenta litri del Decathlon, che l'addetta conferma essere la soluzione ottimale. Decidiamo però di aspettare fino all'imbarco e di passare i controlli tenendo ancora la scatola a protezione dei nostri tesori.

Al metal detector non incontriamo difficoltà, ma i problemi riemergono al controllo passaporti, che avviene automaticamente tramite scansione del documento e riconoscimento facciale. Io passo al primo tentativo, ma Eli non è altrettanto fortunata. Prova una decina di volte senza successo, anche con l'aiuto di un'assistente, che alla fine le indica di recarsi al controllo manuale allo sportello undici. Ricordo bene questo numero perché è l'unica cosa che capiamo di quanto ci dice la ragazza, il cui inglese è piuttosto incerto.

Allo sportello la coda è nuovamente lunghissima e nel frattempo l'orario del volo si avvicina. Mentre i minuti passano inesorabilmente ringrazio il cielo per il ritardo del volo successivo, che ci concede un po' di tempo in più. Finalmente arriva il turno di Elisa e il controllo si svolge rapido e senza intoppi, anche se non posso fare a meno di provare un po' d'invidia: in questo modo si è guadagnata un timbro sul passaporto.

Procediamo velocemente verso il gate d'imbarco, dove gli addetti sono ancora intenti a pranzare, poiché il volo precedente non è ancora arrivato e di conseguenza nemmeno il nostro aereo. Li interrompo con gentilezza e chiedo se sia davvero impossibile portare a bordo la scatola così com'è. Mi confermano quanto già comunicato dall'assistente al check-in e mi chiedono di mostrare il contenuto. Dopo aver verificato che non si tratta di nulla di pericoloso, concordano sul fatto che sia meglio non imbarcarla e ci lasciano a reimpacchettare il tutto.

Vista dall'esterno, la scena deve essere stata memorabile: due passeggeri sudati e agitati che, cercando di mantenere la calma, smantellano una scatola di cartone, estraggono cimeli di terracotta e spargono ritagli di carta di giornale tutt'intorno, per poi infilare il tutto in una borsa leggera che maneggiano come se fosse il Sacro Graal. Il tutto chini per terra nel tempo record di dieci minuti. Il piano procede senza intoppi e la nuova sacca sembra tenere bene, potenzialmente in grado di resistere a urti leggeri.

Soddisfatti del nostro operato, ci prendiamo un momento per tirare il fiato e mangiare qualcosa. Come in ogni viaggio di ritorno che si rispetti, il pranzo è al fast-food e il Burger King ci regala un'enorme abbuffata di panini di pollo croccante e patatine con maionese. Siamo di fretta quindi divoriamo tutto e, dopo esserci dati una ripulita per tornare presentabili, ci dirigiamo di nuovo al gate, dove stanno per iniziare le procedure di imbarco.

La nostra principale preoccupazione è la borsa di tela coi souvenir: vogliamo salire subito per trovarle un posto sicuro e assicurarci che nessuno vi appoggi nulla sopra. Quando saliamo a bordo, ci accorgiamo subito di avere due posti centrali. Accanto a noi ci sono una donna araba completamente vestita di nero con solo gli occhi scoperti, e una donna indonesiana che ha sposato un americano, il quale viaggia con lei insieme ai loro due figli dai bellissimi lineamenti misti. Elisa è quasi sollevata che la donna col burqa sia seduta accanto a lei: se dovesse tossire, almeno non le sputerebbe germi in faccia o nel cibo.

Poco prima della partenza, un ragazzo apre lo scomparto sopra di noi — quello che contiene i cimeli di terracotta. Mi precipito su di lui chiedendo cosa contenesse il sacchetto che aveva appena riposto: si trattava di un paio di sandali e non li aveva nemmeno appoggiati sulla nostra borsa, ma credo di averlo quasi aggredito.

Il cibo a bordo è un capitolo a parte. Stavolta non c'è il rischio che venga contaminato e reso immangiabile, perché peggio di così è impossibile. A Elisa servono riso con manzo. Il riso è terribile — secco e con la carne piena di grasso — quindi quando arriva il mio turno chiedo i noodles: sono finiti, e anch'io mi ritrovo col riso. Insieme ci servono un'insalata di tofu pessima e un budino alla fragola che daresti solo a un gatto anziano che ha perso i denti. L'unica cosa decente che riusciamo a mangiare è il pane, che Elisa riempie con una generosa dose di burro. Siamo grati di esserci fermati al Burger King prima di imbarcarci.

Per noi è notte e durante il volo cerchiamo di dormire qualche ora: stranamente ci riesco anch'io e accumulo un po' di riposo.

\begin{center}
    % Decorative line
    \noindent\rule{0.6\textwidth}{0.4pt}
\end{center}

Quando arriviamo a Jeddah per noi sono quasi le cinque del mattino, ma dobbiamo riportare l'orologio indietro all'una. Qui abbiamo intenzione di mettere in atto un piano di sopravvivenza ben congegnato.

Le alternative erano tre: prenotare una stanza nell'hotel all'interno dell'aeroporto — molto costosa e con disponibilità limitata — oppure prenotare fuori, soluzione di costo medio che comportava però l'ottenimento di visti d'ingresso temporanei, l'uscita dall'aeroporto, avventurarsi per Jeddah di notte in taxi, trovare l'hotel, dormire poco e tornare all'alba per rifare il check-in, il tutto in piena notte e in Arabia Saudita. La terza soluzione, la più economica, consisteva nel dormire in aeroporto nella zona delle partenze, rubando cuscini e coperte dal volo precedente per crearci un giaciglio comodo. Quest'ultima va in porto, tranne per la parte del "comodo."

Ci spostiamo in una zona isolata davanti a un gate vuoto e uniamo due file di sedie per dormire su una piazzola improvvisata. Nonostante le coperte fa un freddo cane e siamo entrambi troppo alti per stare comodi in quel loculo, ma Elisa sembra non curarsi troppo della fisica e riesce a recuperare diverse ore di sonno. Io, su dieci ore di scalo, dormo venti minuti.

Non ci provo nemmeno con troppo impegno: oltre ai problemi già citati c'è una luce fortissima che filtra oltre la mascherina e ogni quarto d'ora un annuncio a tutto volume comunica l'inizio dell'imbarco di qualche volo, prima in arabo e poi in inglese. Come se non bastasse, alle 2:30 l'area si riempie di passeggeri in partenza per il Cairo, e alle quattro il Muezzin — che ormai sembra seguirci come la nuvola di Fantozzi — inizia a pregare a tutto volume dai megafoni. Elisa, però, percepisce a malapena tutti questi fattori disturbanti.

Trascorro quindi il tempo recuperando pagine di diario arretrate e mi rimetto finalmente in pari. Verso le cinque del mattino vado a prendere un cappuccino in un bar e mi viene servito un litro di caffelatte, che era esattamente ciò di cui avevo bisogno.

Comincia ad albeggiare ed Elisa si sveglia. Facciamo due chiacchiere ridendo principalmente dell'assurdità della situazione e poco dopo le sei veniamo allontanati dalla sicurezza insieme ad altri passeggeri perché devono allestire un checkpoint.

Manteniamo la nostra tradizione di cibo da fast-food facendo colazione al McCafé: un americano e un muffin per me, un americano e doppia brioche per l'affamata Elisa. Mangiamo tutto di fronte al nostro gate, che inizia a riempirsi di turisti italiani. Non mi mancava per niente sentire il caos e l'euforia delle scolaresche in gita.

L'imbarco procede senza intoppi e la borsa dei souvenir sale a bordo con noi. Il volo è relativamente breve — solo cinque ore, una bazzecola rispetto al resto — e trascorriamo buona parte del tempo guardando insieme \textit{Batman Begins}, che ci ricorda la grotta di Goa Lawah.

Il pranzo-colazione è più umano: riso con manzo stavolta commestibile, insalata di pomodoro, mais e cetriolo, e una sorta di brownies al cioccolato, oltre alla solita pagnotta salvavita. Con scatoletta di burro per Elisa.

Come all'andata, anche al ritorno Elisa riesce a bloccare il software dello schermo del suo sedile. Nessuna azione criminale da parte sua, ma sembra che questi apparecchi siano particolarmente sensibili, e la sento cliccare sui comandi come se stesse battendo a macchina. Lei ovviamente sostiene che non sia vero, ma non appena prende il mio telecomando — dopo aver bloccato il suo — la luce sopra di me inizia a lampeggiare a intermittenza, segno che sta smanettando anche col mio. Per fortuna manca poco all'arrivo e ne approfitto per continuare a scrivere.

Prima dell'atterraggio l'aereo oscilla un po', ma poi si adagia dolcemente sull'asfalto di Milano.

\begin{center}
    % Decorative line
    \noindent\rule{0.6\textwidth}{0.4pt}
\end{center}

Siamo arrivati in Italia.

Prendiamo i bagagli e saltiamo la lunga coda al controllo passaporti con la scusa "lo scanner non legge i nostri documenti" — così siamo i primi a ritirare i bagagli. Qui arriva però l'amara sorpresa: ci hanno rubato i pacchetti di sigarette che avevamo preso per i nostri amici. Capiamo che anche i loro bagagli sono stati aperti. Seguono attimi di panico in cui Elisa crede che le abbiano rubato anche altro, ma poi trova il bracciale d'argento nella custodia della reflex.

A questo punto dobbiamo solo aspettare il bus per Torino, che però non sembra arrivare. La zona arrivi è un cantiere aperto e non noto la banchina in fondo fino all'ultimo momento. Quando vado a informarmi, l'autista mi dice che se non siamo lì entro un minuto partirà senza di noi.

Sono assonnato, con gambe lente e pesanti, ma scatto attraverso tutta la zona arrivi di Milano Malpensa in direzione di Elisa, mi carico i bagagli e le dico di correre. Con tutto quel peso addosso mi è praticamente impossibile muovermi, ma arriviamo in tempo. Carichiamo e saliamo a bordo. La mia camicia, che aveva resistito eroicamente fino a quel momento, è ora simile a uno straccio bagnato.

Il viaggio verso Torino procede velocemente, fatta eccezione per l'ultima mezz'ora, quando inizia a calare la tensione e la stanchezza prende il sopravvento.

Quando i miei genitori ci accolgono alla stazione di Porta Susa, sono un cadavere. Mi sento come il personaggio di \textit{Se7en} punito per "Accidia" — drogato, denutrito e immobilizzato su un lettino per un anno intero. Ci accompagnano a casa e dopo una doccia mi butto subito a letto, chiedendo a Elisa di svegliarmi per cena, per provare a riadattarmi al fuso orario italiano.

Il risveglio è traumatico, ma mi sforzo di restare sveglio ancora un po' dopo cena per riprendere il ritmo giusto. Il piano funziona, più o meno: verso le cinque del mattino mi sveglio senza riuscire a riaddormentarmi. Mi alzo per andare in bagno, convinto di lavarmi la faccia e iniziare la giornata presto, ma qualcosa va storto. Il solo gesto di alzarmi e cambiare stanza mi fa sentire tutto il peso del viaggio e, appena decido di appoggiarmi di nuovo sul letto, crollo addormentato, vinto dalla stanchezza.


\pagestyle{empty} % disable numbers

\landscapeThinMargin{images/day16/IMG-20240825-WA0009.jpg}

\pagestyle{fancy} % enable numbers